\begin{table}[H]
\centering
\caption{Qualitative Results: Autoregressive Text Generation Samples}
\label{tab:textgen-qualitative}
\begin{tabularx}{\textwidth}{@{}lX@{}}
\toprule
\textbf{Type} & \textbf{Text Sample} \\
\midrule
\textbf{Prompt} & Example prompt in Shona (e.g., "Mangwanani ano...") \\
\textit{Completion} & The model's continuation of the prompt goes here. (e.g., "...tose takamuka zvakanaka.") \\
\midrule
\textbf{Prompt} & Another prompt... \\
\textit{Completion} & Model output... \\
\midrule
\textbf{Prompt} & Complex prompt... \\
\textit{Completion} & Coherent text generation... \\
\bottomrule
\end{tabularx}
\end{table}

\paragraph{Analysis}
The selected samples demonstrate the model's ability to maintain grammatical coherence (subject-verb agreement) and generate contextually relevant continuations. While occasional repetitions occur (common in smaller autoregressive models), the outputs largely adhere to Shona syntax.
